\begin{ex}
Một đoạn thẳng được giới hạn bởi bao nhiêu điểm đầu mút?
\shortans{2}
\loigiai{
Đoạn thẳng bị giới hạn bởi 2 điểm đầu mút.
}
\end{ex}

\begin{ex}
Qua 2 điểm phân biệt, ta vẽ được bao nhiêu đường thẳng?
\shortans{1}
\loigiai{
Qua 2 điểm phân biệt vẽ được duy nhất 1 đường thẳng.
}
\end{ex}

\begin{ex}
Cho 3 điểm A, B, C cùng nằm trên một đường thẳng. Hỏi có bao nhiêu điểm thẳng hàng?
\shortans{3}
\loigiai{
Có 3 điểm thẳng hàng.
}
\end{ex}

\begin{ex}
Cho 3 điểm không thẳng hàng, nối từng cặp điểm ta được bao nhiêu đoạn thẳng?
\shortans{3}
\loigiai{
Nối được 3 đoạn thẳng tạo thành một hình tam giác.
}
\end{ex}

\begin{ex}
Cho 4 điểm trong đó không có 3 điểm nào thẳng hàng. Nối từng cặp điểm ta vẽ được tất cả bao nhiêu đoạn thẳng?
\shortans{6}
\loigiai{
Số đoạn thẳng: $4 \times 3 : 2 = 6$ đoạn thẳng.
}
\end{ex}

\begin{ex}
Trên một đường thẳng lấy 4 điểm phân biệt A, B, C, D theo thứ tự đó. Có bao nhiêu đoạn thẳng trên hình vẽ?
\shortans{6}
\loigiai{
Các đoạn thẳng: AB, AC, AD, BC, BD, CD. Có 6 đoạn thẳng.
}
\end{ex}

\begin{ex}
Đoạn thẳng AB dài 8 cm, điểm C là điểm nằm giữa A và B sao cho AC dài 3 cm. Độ dài đoạn thẳng CB là bao nhiêu xăng-ti-mét?
\shortans{5}
\loigiai{
Độ dài đoạn CB: $8 - 3 = 5$ cm.
}
\end{ex}

\begin{ex}
Cho ba điểm M, N, P thẳng hàng theo thứ tự đó. Biết MN dài 4 cm, NP dài 5 cm. Độ dài đoạn thẳng MP là bao nhiêu xăng-ti-mét?
\shortans{9}
\loigiai{
Độ dài MP: $4 + 5 = 9$ cm.
}
\end{ex}

\begin{ex}
Có bao nhiêu điểm đầu mút trong 3 đoạn thẳng rời nhau?
\shortans{6}
\loigiai{
Mỗi đoạn có 2 đầu mút, 3 đoạn có: $3 \times 2 = 6$ đầu mút.
}
\end{ex}

\begin{ex}
Một hình tam giác có bao nhiêu đỉnh (điểm)?
\shortans{3}
\loigiai{
Hình tam giác có 3 đỉnh.
}
\end{ex}

\begin{ex}
Một hình tứ giác có bao nhiêu đỉnh (điểm)?
\shortans{4}
\loigiai{
Hình tứ giác có 4 đỉnh.
}
\end{ex}

\begin{ex}
Trên đoạn thẳng AB dài 12 cm, người ta đánh dấu thêm điểm O sao cho OA dài 6 cm. Điểm O chia đoạn AB thành 2 đoạn bằng nhau. Độ dài đoạn OB là bao nhiêu xăng-ti-mét?
\shortans{6}
\loigiai{
Độ dài OB: $12 - 6 = 6$ cm.
}
\end{ex}

\begin{ex}
Cho 5 điểm thẳng hàng liên tiếp cách đều nhau 2 cm. Khoảng cách từ điểm thứ nhất đến điểm thứ năm là bao nhiêu xăng-ti-mét?
\shortans{8}
\loigiai{
Có 4 khoảng cách, mỗi khoảng 2 cm: $2 \times 4 = 8$ cm.
}
\end{ex}

\begin{ex}
Có bao nhiêu đường cong khép kín trong hình vẽ gồm 2 hình tròn và 3 hình vuông?
\shortans{2}
\loigiai{
Hình tròn là đường cong khép kín, có 2 hình tròn.
}
\end{ex}

\begin{ex}
Một đường gấp khúc gồm 2 đoạn thẳng nối tiếp nhau có bao nhiêu điểm đỉnh (tính cả 2 đầu mút)?
\shortans{3}
\loigiai{
Có 3 điểm.
}
\end{ex}

\begin{ex}
Một đường gấp khúc gồm 3 đoạn thẳng nối tiếp nhau có bao nhiêu điểm?
\shortans{4}
\loigiai{
Có 4 điểm.
}
\end{ex}

\begin{ex}
Đoạn thẳng CD dài 15 cm. Điểm K nằm giữa C và D sao cho KD dài 7 cm. Đoạn CK dài bao nhiêu xăng-ti-mét?
\shortans{8}
\loigiai{
Đoạn CK dài: $15 - 7 = 8$ cm.
}
\end{ex}

\begin{ex}
Có 3 điểm thẳng hàng X, Y, Z. Biết XY dài 6 cm, YZ dài 8 cm. Hỏi đoạn XZ dài bao nhiêu xăng-ti-mét nếu Y nằm giữa X và Z?
\shortans{14}
\loigiai{
Đoạn XZ dài: $6 + 8 = 14$ cm.
}
\end{ex}

\begin{ex}
Trên một sợi dây thẳng dài 20 cm, người ta thắt 3 nút chia sợi dây thành 4 phần bằng nhau. Mỗi phần dài bao nhiêu xăng-ti-mét?
\shortans{5}
\loigiai{
Mỗi phần dài: $20 : 4 = 5$ cm.
}
\end{ex}

\begin{ex}
Số góc vuông trong một hình chữ nhật là bao nhiêu?
\shortans{4}
\loigiai{
Hình chữ nhật có 4 góc vuông.
}
\end{ex}

\begin{ex}
Trong các chữ cái in hoa: A, B, C, D, O, chữ cái nào được tạo bởi một đường cong khép kín tròn? Chữ O gồm bao nhiêu nét cong khép kín?
\shortans{1}
\loigiai{
Chữ O gồm 1 đường cong khép kín.
}
\end{ex}

\begin{ex}
Một người cắm 4 cọc tiêu thẳng hàng, mỗi cọc cách nhau 3 mét. Khoảng cách từ cọc đầu tiên đến cọc cuối cùng là bao nhiêu mét?
\shortans{9}
\loigiai{
Có 3 khoảng cách: $3 \times 3 = 9$ mét.
}
\end{ex}

\begin{ex}
Để nối 5 điểm thẳng hàng liên tiếp cần bao nhiêu đoạn thẳng nối tiếp nhau?
\shortans{4}
\loigiai{
Cần 4 đoạn thẳng nối tiếp.
}
\end{ex}

\begin{ex}
Một tam giác có chu vi là 18 cm, hai cạnh dài lần lượt là 6 cm và 5 cm. Cạnh thứ ba dài bao nhiêu xăng-ti-mét?
\shortans{7}
\loigiai{
Cạnh thứ ba dài: $18 - 6 - 5 = 7$ cm.
}
\end{ex}

\begin{ex}
Cho đoạn thẳng PQ dài 10 cm. Điểm I là trung điểm của PQ. Đoạn PI dài bao nhiêu xăng-ti-mét?
\shortans{5}
\loigiai{
Đoạn PI dài: $10 : 2 = 5$ cm.
}
\end{ex}
